%% kapa: exam=1 round=30 year=2019 period=1 form=A subject=real_estate subject_ko=부동산학원론 seq=19 number=99
%% source: Q-Net 감정평가사 기출문제 (cst00309) artlSeq=5205205 file=extracted/1008/5205205/2019년 제30회 감정평가사 1차시험 문제지/2019년 제30회 감정평가사 1차시험 문제지/2019년도 제30회 감정평가사 1차 1교시 A형.hwp
%% conversion: kordoc 4.24.51 -> tools/gen1.py (auto); review=proofread-sonnet; flags=none
\begin{question}{99}
지역분석과 개별분석에 관한 설명으로 옳은 것은?
\begin{choices}
  \item 지역분석은 일반적으로 개별분석에 선행하여 행하는 것으로 그 지역 내의 최유효이용을 판정하는 것이다.
  \item 인근지역이란 대상부동산이 속한 지역으로 부동산의 이용이 동질적이고 가치형성요인 중 개별요인을 공유하는 지역이다.
  \item 유사지역이란 대상부동산이 속하지 아니하는 지역으로서 인근지역과 유사한 특성을 갖는 지역이다.
  \item 개별분석이란 지역분석의 결과로 얻어진 정보를 기준으로 대상부동산의 가격을 표준화ㆍ일반화시키는 작업을 말한다.
  \item 지역분석 시에는 균형의 원칙에, 개별분석 시에는 적합의 원칙에 더 유의하여야 한다.
\end{choices}
\end{question}
\answer{3}
